\clearpage
\SourceLeaf{194}
\thispagestyle{empty}
\pdfbookmark[0]{General solution of the problem of mapping parts of a given surface to another}{abbildung}
\vspace*{.35in}
\begin{center}
{\large GENERAL SOLUTION OF THE PROBLEM}\par\bigskip
{\Large OF MAPPING THE PARTS OF A GIVEN SURFACE}\par\bigskip
TO ANOTHER GIVEN SURFACE IN SUCH A WAY\par\bigskip
THAT THE IMAGE IS SIMILAR TO THE OBJECT MAPPED\par\bigskip
IN THEIR SMALLEST PARTS\par\vspace{28pt}
BY\par\vspace{28pt}
{\large CARL FRIEDRICH GAUSS}\par\vspace{30pt}
{\small AS AN ANSWER TO THE PRIZE QUESTION SET BY THE ROYAL SOCIETY OF SCIENCES\par
AT COPENHAGEN FOR MDCCCXXII.}\par\bigskip
\end{center}
\begin{flushright}‘From these the way is paved to greater things.’\end{flushright}
\vfill
\begin{center}
\rule{.8\linewidth}{.3pt}\par\smallskip
Astronomische Abhandlungen edited by H. C. \textsc{Schumacher}.\par
Third fascicle. Altona 1825.\par\smallskip
\rule{.8\linewidth}{.3pt}
\end{center}
\clearpage
\SourceLeaf{196}
\fancyhead[L]{\small GENERAL SOLUTION OF THE PROBLEM}
\pdfbookmark[1]{Preface, December 1822}{abbildung-vorrede}
The author of this treatise believed he had to regard the choice, twice made, of the problem that forms its subject as evidence of the importance the Royal Society attaches to it. This encouraged him to submit to the Society the solution he had found some time before; otherwise, learning of the prize question so late would have deterred him. He regrets that this circumstance forced him to confine himself almost entirely to the essentials and to indicating a few more immediate applications to map projections and higher geodesy. Without the approaching deadline, he would gladly have pursued the development of certain subsidiary matters further and treated the many applications in higher geodesy in detail; he must now reserve this for another time and place.

December 1822.
\begin{center}\rule{.3\linewidth}{.3pt}\end{center}
\clearpage
\SourceLeaf{198}
\begin{center}
{\large GENERAL SOLUTION OF THE PROBLEM}\par\bigskip
OF MAPPING THE PARTS OF A GIVEN SURFACE\par\medskip
TO ANOTHER GIVEN SURFACE IN SUCH A WAY\par\medskip
THAT THE IMAGE IS SIMILAR TO THE OBJECT MAPPED\par\medskip
IN THEIR SMALLEST PARTS.\par\bigskip
\rule{.25\linewidth}{.3pt}
\end{center}
\renewcommand{\Article}[1]{\par\medskip\Needspace{4\baselineskip}\pdfbookmark[1]{Art. #1}{abbildung-art-#1}\centerline{\textbf{#1.}}\nobreak\smallskip}
\Article{1}
The nature of a curved surface is determined by an equation between the coordinates \(x,y,z\) of each of its points. By virtue of this equation, each of these three variable quantities can be regarded as a function of the other two. It is more general still to introduce two further variable quantities \(t,u\) and represent each of \(x,y,z\) as a function of \(t\) and \(u\); then, at least generally speaking, definite values of \(t\) and \(u\) always correspond to a definite point on the surface, and conversely.

\Article{2}
For a second curved surface, let \(X,Y,Z,T,U\) have meanings analogous to those of \(x,y,z,t,u\), respectively, for the first.

\Article{3}
To \textit{map} the first surface to the second means to specify a rule by which a definite point on the second surface is to correspond to each point on the first. This is done by setting \(T\) and \(U\) equal to definite functions of the two variable quantities \(t\) and \(u\).
\SourceLeaf{199}
If the mapping is to satisfy certain conditions, these functions can no longer be arbitrary. Since \(X,Y,Z\) thereby also become functions of \(t\) and \(u\), these functions must satisfy both the condition prescribed by the nature of the second surface and the condition that the mapping is to fulfill.

